\documentclass[aip, adv, reprint, superscriptaddress]{revtex4-2}
\usepackage{amsmath, amssymb, amsfonts}
\usepackage{graphicx}
\usepackage{hyperref}

\begin{document}

\title{Photonic Quantum Error Correction from Geometric Stability Principles in SEFI/GWFM}

\author{Jason Dutton}
\affiliation{Castlewood, Virginia, USA}

\begin{abstract}
Photonic quantum information platforms are highly susceptible to loss, dephasing, mode 
mixing, and spectral drift. Conventional photonic quantum error correction (QEC) 
strategies provide effective code constructions, but they do not yet rest on a single 
physical principle that explains why errors arise and how correction should be structured.

Here, we introduce a geometric framework for photonic QEC derived from the Single 
Entity Field Interpretation (SEFI) and the Geometric Worldline Field Model (GWFM). In 
this model, photonic errors correspond to geometric displacements on stability surfaces, 
and correction corresponds to restoring the sovereignty of the encoded field. We derive 
stabilizer structures, geometric syndrome relations, and error-surface topology for 
photonic modes, and we outline experimental implementations based on integrated 
photonics and weak or quantum-nondemolition (QND) measurement. We further show 
how SEFI/GWFM offers a geometric interpretation of SPDC-based quantum randomness 
production and Bloch-sphere CHSH Bell-inequality behavior. This framework provides a 
physically grounded foundation for photonic QEC and suggests new bosonic-code designs 
with intrinsic stability advantages.
\end{abstract}

\keywords{photonic quantum error correction, geometric stability, bosonic codes, integrated photonics, quantum information}

\maketitle

\section{Introduction}

Photonic quantum information processing has emerged as a leading platform for quantum 
communication, sensing, and scalable quantum computing. Photons offer long coherence 
times, room-temperature operation, and compatibility with integrated photonics. However, 
photonic systems are susceptible to distinct error channels, including loss, mode 
deformation, spectral wandering, and imperfect interference.

Existing photonic QEC approaches---including cat codes, binomial codes, and GKP 
codes---provide useful mathematical tools for error suppression, but they do not arise 
from a single physical principle that governs error formation or dictates the architecture 
of correction.

SEFI \cite{DuttonSEFI} and GWFM \cite{DuttonGWFM} provide such a principle. SEFI 
formalizes layered field identity through FIELD ORIGIN, FIELD AUTHORSHIP, FIELD 
SOVEREIGNTY, and WARP EXPRESSION. GWFM provides the geometric worldline 
foundation for field propagation, curvature, torsion, and stability.

This paper develops a photonic QEC model grounded in these geometric principles and 
shows that SEFI/GWFM also provides geometric interpretations of SPDC-based QRNG 
and CHSH Bell-inequality behavior.

\section{SEFI/GWFM Stability Geometry for Photonic Modes}

\subsection{Field Origin and Photonic Identity}

In SEFI, FIELD ORIGIN defines the primitive geometric identity of a field. For a photonic 
mode, the origin layer determines the baseline worldline geometry, spectral identity, 
polarization geometry, and curvature constraints. Stability is treated as a primitive 
property through the bound
\begin{equation}
\left\| \frac{d^2 r}{dt^2} \right\| \le S_0,
\end{equation}
where $S_0$ is the sovereignty-bound acceleration limit \cite{DuttonSEFI}.

\subsection{Field Sovereignty and Stability Surfaces}

FIELD SOVEREIGNTY ensures invariance under allowed geometric transformations:
\begin{equation}
F(T(r(t))) = F(r(t)).
\end{equation}
This defines a stability surface: a geometric manifold representing the allowed 
transformations of a photonic mode.

Errors correspond to displacements orthogonal to this surface.

\subsection{GWFM Worldline Geometry}

GWFM defines curvature as
\begin{equation}
K(t) = \frac{\| r'(t) \times r''(t) \|}{\| r'(t) \|^3},
\end{equation}
and treats propagation as a geometric worldline subject to curvature and torsion 
constraints.

Loss corresponds to worldline termination; dephasing corresponds to curvature-induced 
torsion; mode mixing corresponds to displacement into neighboring stability surfaces.

\section{Geometric Photonic Error Surfaces}

\subsection{Loss Errors}

Loss is modeled as collapse of the worldline into a boundary region of the stability 
surface.

\subsection{Dephasing Errors}

Dephasing corresponds to geometric torsion of the stability surface.

\subsection{Mode-Mixing Errors}

Mode mixing corresponds to displacement into a neighboring stability surface.

\subsection{Spectral Wandering}

Spectral drift corresponds to translation along the spectral axis of the stability manifold.

\section{Geometric Stabilizers and Syndrome Extraction}

\subsection{Stabilizer Geometry}

A stabilizer is defined as a geometric constraint that preserves the sovereignty of the 
encoded field. For photonic modes, stabilizers correspond to curvature constraints, 
torsion constraints, spectral-position constraints, and polarization-surface constraints.

\subsection{Syndrome Geometry}

Weak or QND photonic measurements can sample curvature deviation, torsion deviation, 
spectral displacement, and polarization rotation. These measurements define a 
geometric syndrome vector.

\section{SEFI/GWFM Interpretation of SPDC-Based QRNG}

Spontaneous parametric down-conversion (SPDC) produces photon pairs with strong 
temporal and spatial correlations due to energy and momentum conservation 
\cite{HongMandel, Uren2005}. In non-collinear, degenerate SPDC, the photons appear 
on an annular ring whose geometry reflects the underlying momentum-matching 
conditions \cite{Jabir2017}. Recent work has shown that sectioning this ring enables 
high-bit-rate quantum random number generation (QRNG) by exploiting correlations 
between diametrically opposite regions \cite{SPDCQRNG2023}.

Within SEFI, the annular spatial distribution corresponds to a stability surface defined 
by FIELD ORIGIN and FIELD SOVEREIGNTY, where diametrically opposite photon 
positions represent sovereign geometric symmetries of the pair field \cite{DuttonSEFI}. 
The primitive stability constraint ensures that the worldline curvature of each photon 
remains bounded, producing the highly regular annular geometry observed 
experimentally.

GWFM models the propagation of each photon as a geometric worldline constrained by 
curvature and torsion, providing a physical basis for the strong temporal correlations 
observed between spatially separated sections of the SPDC ring \cite{DuttonGWFM}. 
Randomness extraction corresponds to sampling warp-expression variations across 
distinct worldline bundles, enabling multi-bit QRNG without beam-splitter-based 
probabilistic routing.

\section{SEFI/GWFM Interpretation of CHSH Bell-Inequality Geometry}

Recent experimental work has demonstrated CHSH Bell-inequality violations across the 
full Poincar\'e--Bloch sphere, including both Bell and non-Bell maximally entangled 
states \cite{CardosoIsidoro2026}. These results show that CHSH violation is strongly 
dependent on the geometric relationship between the polarization bases used for each 
photon.

Within SEFI, polarization-entangled photon pairs occupy distinct regions of a stability 
surface defined by FIELD ORIGIN and FIELD SOVEREIGNTY \cite{DuttonSEFI}. 
Sovereignty requires that allowed geometric transformations preserve field identity,
\begin{equation}
F(T(r(t))) = F(r(t)),
\end{equation}
which implies that CHSH violation occurs only when measurement bases correspond to 
transformations lying within the sovereignty-preserving submanifold of the stability 
surface.

GWFM provides a complementary interpretation: polarization evolution is treated as a 
worldline on the Poincar\'e--Bloch sphere, with curvature and torsion determining the 
geometric relationship between measurement bases. Bell states correspond to worldlines 
with minimal torsion, enabling CHSH violation across a wide range of basis choices, 
whereas non-Bell states correspond to worldlines with higher torsion, producing 
basis-dependent suppression of violation.

\section{Experimental Implementation}

\subsection{Integrated Photonics}

Waveguide-based photonic circuits can implement geometric stabilizers, weak or QND 
sampling, spectral-geometry control, and mode-geometry correction.

\subsection{Weak/QND Measurement}

QND detectors provide geometric sampling without worldline termination.

\subsection{Photonic Syndrome Extraction}

Syndromes are extracted through homodyne sampling, spectral interferometry, 
polarization tomography, and mode-shape reconstruction.

\subsection{Correction Operations}

Correction corresponds to geometric restoration: spectral translation, phase-geometry 
correction, mode-shape reformation, and polarization-surface restoration.

\section{Conclusion}

We have introduced a geometric framework for photonic quantum error correction derived 
from SEFI/GWFM stability principles. Photonic errors correspond to geometric 
displacements on stability surfaces, and correction corresponds to restoring field 
sovereignty. We further show that SEFI/GWFM provides geometric interpretations of 
SPDC-based QRNG and CHSH Bell-inequality behavior. This establishes a physically 
grounded foundation for photonic QEC and opens a new direction for geometric quantum 
information theory.

\section*{Author Declarations}

\subsection*{Conflict of Interest}
The authors have no conflicts to disclose.

\section*{Data Availability}
No datasets were generated or analyzed during the current study.

\bibliographystyle{aipnum4-2}

\begin{thebibliography}{99}

\bibitem{DuttonSEFI}
J. Dutton, \textit{SEFI: The Single Entity Field Interpretation --- Formal Mathematical Construction}, 
Manuscript JMP26-AR-02501, submitted to the Journal of Mathematical Physics (2026).

\bibitem{DuttonGWFM}
J. Dutton, \textit{Geometric Worldline Foundations of the Electron Field}, 
Manuscript JMP26-AR-02470, submitted to the Journal of Mathematical Physics (2026).

\bibitem{HongMandel}
C. K. Hong, Z. Y. Ou, and L. Mandel, Phys. Rev. Lett. 59, 2044 (1987).

\bibitem{Uren2005}
A. B. U'Ren et al., Phys. Rev. A 72, 021802 (2005).

\bibitem{Jabir2017}
M. V. Jabir and G. K. Samanta, Sci. Rep. 7, 12613 (2017).

\bibitem{SPDCQRNG2023}
S. Singh et al., Adv. Quantum Technol. 6, 2300177 (2023).

\bibitem{CardosoIsidoro2026}
C. Cardoso-Isidoro and E. J. Galvez, 
\textit{Clauser–Horne–Shimony–Holt Bell-inequality violability with the full Poincaré–Bloch sphere}, 
AVS Quantum Sci. 8, 021401 (2026).

\end{thebibliography}

\end{document}
